% Generated draft: requires mathematical and source review.
\documentclass[11pt]{article}
\usepackage[margin=25mm]{geometry}
\usepackage{fontspec}
\setmainfont{Latin Modern Roman}
\usepackage{amsmath,amssymb,mathtools,unicode-math}
\setmathfont{Latin Modern Math}
\usepackage{enumitem,needspace,xurl,hyperref}
\setlength{\parindent}{0pt}
\setlength{\parskip}{5pt}
\setlength{\emergencystretch}{3em}
\newcommand{\N}{\mathbb{N}}
\newcommand{\Z}{\mathbb{Z}}
\newcommand{\Q}{\mathbb{Q}}
\newcommand{\R}{\mathbb{R}}
\newcommand{\C}{\mathbb{C}}
\newcommand{\F}{\mathbb{F}}
\newcommand{\A}{\mathbb{A}}
\newcommand{\PP}{\mathbb{P}}
\newcommand{\E}{\mathbb{E}}
\newcommand{\eps}{\varepsilon}
\newcommand{\dd}{\,\mathrm{d}}
\begin{document}
\section*{UnsolvedMath ID: 20001156}
\textbf{Source problem code:} catalogue entry
\subsection{Definitions and mathematical statement}
In this source a finite simple graph is balanced if every induced cycle has length divisible by four; in particular it is bipartite. The word ``basic'' is specific to this balanced-graph question: there is a bipartition with every vertex on one side having degree at most two, or there is a hole $Q$ such that $G[V(G)\setminus V(Q)]$ is complete bipartite. This is not the basic-class list for general Berge decompositions.

A skew-partition split here is a partition into four nonempty sets $A,B,C,D$ with $A$ complete to $B$ and $C$ anticomplete to $D$. The cross-edge structure of a $2$-join is a partition $X_1,X_2$ with disjoint nonempty attachments $A_i,B_i\subseteq X_i$, where cross edges are exactly the complete $A_1$--$A_2$ and $B_1$--$B_2$ pairs. Usual definitions add size and side-nondegeneracy conditions; the exact proper-$2$-join version used by the attached aid is not supplied and must be checked, rather than silently inferred from the cross-edge structure alone.

The Conforti--Cornu\'ejols--Rao conjecture in the source is
\[
 \forall G:\quad G\text{ balanced}\ \Longrightarrow\
 \bigl[G\text{ basic}\ \text{or }G\text{ has a $2$-join}\
 \text{or }G\text{ has a skew-partition}\bigr].
\]
No connectedness assumption is present. The aid's established special cases retain that scope: the conclusion, with a proper $2$-join or a four-nonempty-part skew split, holds whenever the balanced graph is induced-$C_4$-free or has maximum degree at most three.

The aid describes deriving the connected case from a published balanceable-graph decomposition theorem, excluding $6$-joins and its named block $R_{10}$ using induced six-cycles, converting remaining star cutsets to skew partitions except for the basic graph $P_3$, and extending to disconnected graphs. A star cutset is a separating vertex set containing a vertex adjacent to all other vertices of the set. The definitions and theorem details for balanceability, $6$-joins, $R_{10}$, and proper $2$-joins are absent. The unrestricted conjecture is explicitly not resolved by the aid, whose checked literature did not establish its present status.
\subsection{Short English statement}
Ask whether every balanced graph is basic or admits a $2$-join or skew partition. The supplied aid gives the exact trichotomy for induced-four-cycle-free and subcubic balanced graphs, including disconnected graphs, while leaving the general conjecture and some decomposition conventions unresolved.
\subsection{Sources}
\begin{itemize}
\item AIM workshop \emph{The Perfect Graph Conjecture}, Chapter K, ``Balanced Graphs,'' Conjecture 1 attributed to Conforti--Cornu\'ejols--Rao: \url{https://aimath.org/WWN/perfectgraph/perfectgraph.pdf}.
\item Local registry \texttt{data/unsolvedmath/problems.json}, permanent ID 20001156, exact source code \texttt{catalogue entry}; canonical location \texttt{aim-combinatorics-notes.json notes[280]}. Companion ID 20001155 (\texttt{AIM-COMBINATORICS-0280}, \texttt{notes[279]}) ends with Chapter K's explicit balanced/basic definitions.
\item Attached attempt 1, partial result, problem status at run unknown, rigor proved, low novelty confidence. It invokes a published balanceable-graph decomposition theorem without a full citation; it is machine-generated, not peer reviewed, requiring independent verification before citation.
\end{itemize}
\end{document}
